

\section{Phase 1: Proposing Ideas for Multimodal Fusion - Done}
\begin{longtable}{|p{8cm}|p{7cm}|}
\hline
\textbf{Work} & \textbf{In charge} \\
\hline
\endhead
Research on proposed pipeline
& Nhat Khang, Tri Duc\\
\hline
Research on datasets for training
& Nhat Khang, Tri Duc\\
\hline
Research on CLIP model and related works
& Tri Duc \\
\hline
Theoretical knowledge preparation for the project
& Nhat Khang\\
\hline
\end{longtable}

\section{Phase 2: Re-adjusting Thesis Topic and Work On New Method for Multimodal Fusion}
\label{sec:phase2}

\newpage
\clearpage % Clear any previous content


\thispagestyle{empty}
\begin{landscape} % Start the landscape environment
    \begin{figure}[H]
        \centering
        \includegraphics[width=1.7\textheight]{sections/appendixes/plan-khang.png}
    \end{figure}
\end{landscape} % End the landscape environment


% \noindent In Phase 2 of our project, we will focus on several key tasks aimed at enhancing the capabilities and performance of our multimodal inference system. Specifically, we will explore methods for hybridizing information derived from both text embedding and image embedding within the CLIP framework. This research aims to enhance the synergy between textual and visual information in our multimodal inference system, ultimately leading to more nuanced and contextually rich responses. By delving into the potential of information encoding from both modalities, we anticipate further advancements in the system's ability to comprehend and generate responses based on diverse multimodal inputs.

% \noindent The design and implementation of a Vector Database will be crucial for efficiently storing and managing the vector representations of images and associated metadata. This database will serve as the backbone of our system, enabling quick and effective retrieval of relevant information during the inferencing process.

% \noindent Next, the implementation of the RAG model will be a pivotal step in our project. RAG integrates information retrieval with text generation, allowing for more contextually relevant and coherent responses. Connecting PhoCLIP, our pretrained CLIP model fine-tuned for the Vietnamese language, with RAG will further enhance the system's ability to understand and generate responses based on multimodal inputs, thereby enabling responses with both text and images.

% \noindent Additionally, we will embark on the task of crawling fashion data from various fashion websites. This data will be used for further fine-tuning and testing of the RAG system and the entire pipeline. By incorporating real-world fashion data into our system, we aim to improve its accuracy and relevance in providing fashion-related recommendations and assistance.

% \noindent Once the system components are in place, extensive system testing will be conducted to ensure functionality, reliability, and performance under various scenarios and conditions. This testing phase will involve rigorous evaluation of individual components as well as the integrated system as a whole. Comprehensive testing protocols will be devised to assess the system's robustness and accuracy across a diverse range of inputs and use cases.

% \noindent Finally, the overall results of the system will be evaluated based on predetermined metrics and criteria. These evaluations will provide insights into the system's effectiveness, efficiency, and suitability for real-world applications, guiding further iterations and refinements to optimize performance and user satisfaction.